%-------------------------------------------------------------
% Template para apresentações
% WebMedia 2022
% webmedia.org.br
%-------------------------------------------------------------

\documentclass{beamer}
\usepackage[portuges]{babel}
\usepackage{hyperref}
\usepackage[T1]{fontenc}
\usepackage{tikz}
\usepackage{listings,textcomp,color}
\usepackage{latexsym,xcolor,multicol,booktabs,calligra}
\usepackage{amsmath,amssymb,BOONDOX-cal,bm}	
\usepackage{graphicx,pstricks,stackengine}    

% set colors
\definecolor{webmediaTextColor}{RGB}{51,51,51}
\definecolor{webmediaColor}{RGB}{133,42,133}

\setbeamercolor*{palette primary}{bg=webmediaColor}
\setbeamercolor*{palette secondary}{bg=webmediaTextColor, fg = white}
\setbeamercolor*{palette tertiary}{bg=webmediaTextColor, fg = white}
\setbeamercolor*{titlelike}{fg=webmediaTextColor}
\setbeamercolor*{title}{bg=webmediaTextColor, fg = white}
\setbeamercolor*{item}{fg=webmediaTextColor}
\setbeamercolor*{caption name}{fg=webmediaTextColor}
\setbeamercolor{footline}{bg=webmediaTextColor, fg = white}
\setbeamercolor{block title}{bg=webmediaColor, fg=white}

%\titlegraphic{%
%  \begin{picture}(0,0)
%    \put(170,100){\makebox(0,0)[rt]{\includegraphics[width=5cm]%{webmedia22.png}}}
  %\end{picture}}

\titlegraphic{ 
\begin{tikzpicture}[overlay,remember picture]
\node[right=-0.15cm] at (current page.150){
    \includegraphics[width=1.2\textwidth]{logoWebmedia.png}
    \includegraphics[width=1.45\textwidth]{webmedia22.png}
};
\end{tikzpicture}
}

\usebackgroundtemplate{%
\tikz[overlay,remember picture] \node[opacity=0.1, at=(current page.center)] {
   \includegraphics[height=\paperheight,width=\paperwidth]{mon}};
}

\title{Exemplo de Título}
\subtitle{Template de Apresentações do WebMedia 2022}
\author[Tal1 e Tal2]{Fulana de Tal1\inst{1} \and Ciclano de Tal2\inst{2}}
\date[Nov'22]{11 de Novembro de 2022}
\institute{
  \inst{1} Departamento de Informática\\
  Universidade Tecnológica Federal do Paraná\and
  \inst{2}Departamento de Matemática\\
  Universidade Federal de XX
}

\usepackage{webmedia}

\def\cmd#1{\texttt{\color{red}\footnotesize $\backslash$#1}}
\def\env#1{\texttt{\color{blue}\footnotesize #1}}

\lstset{language=Python,upquote=true,
  basicstyle=\ttfamily\footnotesize,numbers=left,
  numberstyle=\tiny,stepnumber=1,numbersep=5pt,
  backgroundcolor=\color{white},frame=single,tabsize=2,
  showspaces=false,showstringspaces=false,showtabs=false,
  breaklines=true,breakatwhitespace=true,escapeinside=||,
  commentstyle=\color{black!80}\it
}


\newtheorem{thm}{Theorem}[theorem]

%—-------------------------------------------------------------
\begin{document}
\begin{frame}
\titlepage
\end{frame}
\usebackgroundtemplate{ } 

\begin{frame}{Sumário}
\tableofcontents[sectionstyle=show,subsectionstyle=show/shaded/hide,subsubsectionstyle=show/shaded/hide]

Para sumário automático, inclua seções.
\end{frame}

%—------------------------------------------------------
%\section{Introdução}
	    
\begin{frame}{Introdução}
Lorem ipsum dolor sit amet, consectetur adipiscing elit, sed do eiusmod tempor incididunt ut labore et dolore magna aliqua:

\vskip 0.25cm
\begin{block}{Lorem ipsum}
Ut enim ad minim veniam, quis nostrud exercitation ullamco laboris nisi ut aliquip ex ea commodo consequat. 
\end{block}
\end{frame}
	    
\begin{frame}{Objetivos}
Lorem ipsum dolor sit amet, consectetur adipiscing elit, sed do eiusmod tempor incididunt ut labore et dolore magna aliqua:
\begin{itemize} 
\item Ut enim ad minim veniam, quis nostrud exercitation ullamco laboris nisi ut aliquip ex ea commodo consequat. 
\item Ut enim ad minim veniam, quis nostrud exercitation ullamco laboris nisi ut aliquip ex ea commodo consequat. 
\end{itemize}
\end{frame}

%—------------------------------------------------------
%\section{Método}
	    
\begin{frame}{Metodologia}
\begin{itemize} 
\item ...
\end{itemize}
\end{frame}

\begin{frame}{Metodologia}
\begin{columns}[T]
    \begin{column}<0->{.5\textwidth}
    \begin{figure}[thpb]
    \centering
    \resizebox{1\linewidth}{!}{\includegraphics{mon.jpg}}
    \caption{Test}
    \label{fig:1}
    \end{figure}
    \end{column}%
    \hfill%
    \begin{column}<0->{.5\textwidth}
    \begin{itemize}
        \item<1-> Blah blah blah
	\begin{itemize}
	   \item<1-> Blah blah blah
	\end{itemize}
	\item<2-> Blah blah blah
	\begin{itemize}
	    \item<2-> Blah blah blah
	\end{itemize}
    \end{itemize}
\end{column}
\end{columns}
\end{frame}


\begin{frame}[fragile]{Exemplo de código}

Exemplo de código em Python.

\begin{lstlisting}[language=Python,keywordstyle=\bf,stringstyle=\it]
# Fibonacci numbers module

def fib(n):    # write Fibonacci series up to n
    a, b = 0, 1
    while a < n:
        print(a, end=' ')
        a, b = b, a+b
    print()
\end{lstlisting}
  
\begin{block}{Exemplo de bloco}
Ut enim ad minim veniam, quis nostrud exercitation ullamco laboris nisi ut aliquip.
\end{block}
  
\end{frame}
%—------------------------------------------------------
%\section{Resultados}
	    
\begin{frame}{Resultados}
\begin{minipage}{0.47\textwidth}
    \begin{itemize}
        \item First item
        \item Second item
        \item Third item
    \end{itemize}
\end{minipage}
\begin{minipage}{0.5\textwidth}
    \rule{\textwidth}{0.75\textwidth}
\end{minipage}
\end{frame} 

\begin{frame}{Resultados}
Lorem ipsum dolor sit amet, consectetur adipiscing elit, sed do eiusmod tempor incididunt ut labore et dolore magna aliqua:

\vskip 0.25cm 
\begin{table}[!h]
\centering
\begin{tabular}{l|cccc}
& Mean & Sd. Dev. & 95\% HPD \\ \hline
$\mu_1$ & 1.220 & 0.303 & [0.567, 1.821] \pause\\
$\mu_2$ & 2.676 & 0.409 & [1.863, 3.498] \pause\\
$\rho$ & 0.313 & 0.264 & [-0.295, 0.749]
\end{tabular}
\end{table}
\end{frame} 

%—------------------------------------------------------
%\section{Conclusões}
	    
\begin{frame}{Conclusões}
Lorem ipsum dolor sit amet, consectetur adipiscing elit, sed do eiusmod tempor incididunt ut labore et dolore magna aliqua.
\end{frame}


\nocite{*}
\begin{frame}[allowframebreaks]
\frametitle{Bibliography}
\begin{thebibliography}{10}
\beamertemplatebookbibitems

\bibitem{Tankov} Peter Tankov.
\newblock {\em Financial Modelling with Jump Processes}. \newblock Editora Taylor \& Francis , ano $2003$.

\bibitem{Sato} Sato, K.
\newblock {\em Levy Processes and Infinitely Divisible Distributions}. \newblock Editora Cambridge University Press , ano $1999$.

\beamertemplatearticlebibitems

\bibitem{Protter} Protter, P.
\newblock {\em Stochastic Integration and Differential Equations}. \newblock U.S. Government Printing Office , ano $2004$.

\end{thebibliography}
\end{frame}

\end{document}
